\noindent\textbf{Elección del procedimiento.} Aplicar $a=\partial_tu+u\partial_xu$ manteniendo fija la posición al derivar respecto al tiempo. Aquí $\partial_tu=\beta(1+x/L)$ y $\partial_xu=U(t)/L$.

\begin{enumerate}

\item Calcular la amplitud del campo en el instante dado.

\[\fitmath{\mathrm{Ut}=U_{0}+\beta\cdot t}\]

\[\fitmath{\mathrm{Ut}=(10)+(5)\cdot (1)\approx 15\,\mathrm{m/s}}\]

\item Introducir la posición para la velocidad local.

\[\fitmath{\mathrm{uloc}=\mathrm{Ut}\cdot \left(1+\frac{x}{L}\right)}\]

\[\fitmath{\mathrm{uloc}=(15)\cdot \left(1+\frac{(0.25)}{(0.5)}\right)\approx 22.5\,\mathrm{m/s}}\]

\item La derivada temporal mantiene $x$ fijo.

\[\fitmath{\mathrm{aloc}=\beta\cdot \left(1+\frac{x}{L}\right)}\]

\[\fitmath{\mathrm{aloc}=(5)\cdot \left(1+\frac{(0.25)}{(0.5)}\right)\approx 7.5\,\mathrm{m/s^{2}}}\]

\item Multiplicar la velocidad local por su gradiente.

\[\fitmath{\mathrm{aconv}=\frac{\mathrm{uloc}\cdot \mathrm{Ut}}{L}}\]

\[\fitmath{\mathrm{aconv}=\frac{(22.5)\cdot (15)}{(0.5)}\approx 675\,\mathrm{m/s^{2}}}\]

\item Ambas aceleraciones apuntan en $x$ y se suman con signo.

\[\fitmath{\mathrm{atot}=\mathrm{aloc}+\mathrm{aconv}}\]

\[\fitmath{\mathrm{atot}=(7.5)+(675)\approx 682.5\,\mathrm{m/s^{2}}}\]

\end{enumerate}
